% !TEX root = ../main.tex
% !TEX program = lualatex
\documentclass[../main.tex]{subfiles}
\IfSubfilesClassLoaded{\externaldocument{\subfix{../build/main}}}{}
% ====================
% File: 15_policy.tex
% ====================
\begin{document}
\IfSubfilesClassLoaded{\chapter{EDO Study Guide}}{}
% === POLICY (EDO 3.3.1) ===
\section{Acquisition Policy and Players}

\subsection{Summary}
\begin{description}
  \item[\textbf{Rules and regulations.}] Acquisition authority flows from acquisition-related law, Federal regulations (\ac{far}, \ac{dfars}, and component supplements), and \ac{dod}/\ac{don} policy documents in the 5000 series and \ac{secnavinst} guidance~\autocite{edo-3-3-1-acq-policy-players-2025}.
  \item[\textbf{Decision support systems.}] \ac{jcids} defines requirements, \ac{ppbe} provides resources, and the \ac{das} delivers materiel; effective program management aligns all three~\autocite{edo-3-3-1-acq-policy-players-2025}.
  \item[\textbf{WAS emphasis.}] \ac{dod} redesignated the \ac{das} as the \ac{was} to bias acquisition toward speed-to-capability, commercial-first approaches, and portfolio-level outcomes over process compliance~\autocite{secwar-was-memo-2025,secwar-acq-transformation-2025}.
  \item[\textbf{Risk and affordability.}] Cost, schedule, and performance risks define the trade space, while affordability keeps life-cycle cost within disciplined bounds~\autocite{edo-3-3-1-acq-policy-players-2025}.
  \item[\textbf{Key acquisition players.}] Acquisition Program, \ac{dae}, \ac{cae}, \ac{mdauthority}, \ac{peo}, and \ac{pm} definitions establish who owns decisions and reporting~\autocite{edo-3-3-1-acq-policy-players-2025}.
  \item[\textbf{Oversight and reporting.}] \ac{acat} categories scale oversight by cost and special interest; reporting relies on \ac{rdais}, \ac{daes}, \ac{sar}, \ac{ucr}, and \ac{dave}~\autocite{edo-3-3-1-acq-policy-players-2025}.
\end{description}

\subsection{Rules and Regulations}
\begin{enumerate}
  \item Acquisition-related laws.
  \item Federal Government regulations.
  \item \ac{dod} and \ac{don} acquisition policy documents~\autocite{edo-3-3-1-acq-policy-players-2025}.
\end{enumerate}

\subsection{Acquisition-Related Law}
Statutory authority from Congress provides the legal basis for systems acquisition. Title~10, U.S.C.\ governs the organization, structure, and operation of the Armed Forces, and the annual \ac{ndaa} is a primary mechanism for modifying acquisition authorities and reporting requirements~\autocite{edo-3-3-1-acq-policy-players-2025}.

\begin{itemize}
  \item Armed Services Procurement Act (1947), as amended.
  \item Small Business Act (1963), as amended.
  \item Office of Federal Procurement Policy Act (1983), as amended.
  \item Competition in Contracting Act (1984).
  \item Department of War Procurement Reform Act (1985).
  \item Department of War Reorganization Act (Goldwater-Nichols) (1986).
  \item Defense Acquisition Workforce Improvement Act (1990).
  \item Federal Acquisition Streamlining Act (1994).
  \item Clinger-Cohen Act (1996).
  \item Weapon Systems Acquisition Reform Act (2009)~\autocite{edo-3-3-1-acq-policy-players-2025}.
\end{itemize}

\subsection{Federal Government Regulations}
\begin{itemize}
  \item The \ac{far} is the primary regulation for all Federal agencies acquiring supplies and services with appropriated funds; it guides acquisition planning, competition, contract award, and warranty policy.
  \item The \ac{dfars} supplements the \ac{far} with \ac{dod}-specific procedures.
  \item Component-unique supplements (for example, \ac{nmcars} for the Navy) tailor policy at the Service level.
\end{itemize}
The \ac{far} is the highest-level acquisition regulation~\autocite{edo-3-3-1-acq-policy-players-2025}.

\subsection{Revolutionary FAR Overhaul Class Deviations}
Presidential action established FAR-overhaul policy direction through ``Restoring Common Sense to Federal Procurement''~\autocite{WH-EO-Restoring-Common-Sense-Procurement-2025}. OMB Memo M-25-26 provides implementation direction for agencies and the \ac{far} Council during execution~\autocite{OMB-M-25-26-Overhauling-FAR-2025}. FAR Council deviation guidance then supports interim operations while part-by-part text is revised and final rulemaking proceeds~\autocite{FAR-Council-Deviation-Guidance-RFO-2025}. DoD class deviations posted by \ac{dpcap} currently cover Part~6 (competition), Part~10 (market research), Part~12 (commercial products and services), and Part~18 (emergency acquisitions), effective Feb~1,~2026~\autocite{DoD-RFO-Deviation-Part6-2026,DoD-RFO-Deviation-Part10-2026,DoD-RFO-Deviation-Part12-2026,DoD-RFO-Deviation-Part18-2026}. The Acquisition.gov deviation guide is the live part-by-part status tracker for agency implementation~\autocite{FAR-Overhaul-Deviation-Guide-2026,FAR-Overhaul-RFO-2026}. For source authority, use the EO, OMB memo, and FAR Council guidance as controlling implementation references; White House fact-sheet and DAU materials are informative context~\autocite{WH-FactSheet-Restoring-Common-Sense-Procurement-2025,DAU-FAR-Overhaul-Announcement-2025}.

\subsection{Acquisition Policy Documents}
\begin{itemize}
  \item \ac{dodd}~5000.01 (Defense Acquisition System) sets policy and assigns senior acquisition responsibilities.
  \item \ac{dodi}~5000.02 (Operation of the Adaptive Acquisition Framework) establishes six pathways to tailor acquisition strategies.
  \item Acquisition Guidebooks provide non-mandatory guidance, lessons learned, and best practices.
  \item \ac{secnavinst}~5000.2 and the SECNAV M~5000.2 DoN Acquisition and Capabilities Guidebook provide \ac{don}-specific policy and tailoring guidance~\autocite{edo-3-3-1-acq-policy-players-2025}.
\end{itemize}

\subsection{Adaptive Acquisition Framework}
The \ac{aaf} supports the National Defense Strategy by providing tailored pathways for delivering capability. Pathway instructions include:
\begin{itemize}
  \item Major Capability Acquisition (\ac{mca}) -- \ac{dodi}~5000.85.
  \item Urgent Capability Acquisition -- \ac{dodi}~5000.81.
  \item Middle Tier of Acquisition (\ac{mta}) -- \ac{dodi}~5000.80.
  \item Software Acquisition -- \ac{dodi}~5000.87.
  \item Defense Business Systems -- \ac{dodi}~5000.75.
  \item Defense Acquisition of Services -- \ac{dodi}~5000.74.
\end{itemize}
Functional area guidance includes engineering, test and evaluation, cybersecurity, \ac{aoa}, cost estimating, intellectual property, protection, human systems integration, intelligence, and information technology~\autocite{edo-3-3-1-acq-policy-players-2025}.

\begin{figure}[H]
  \centering
  \includegraphics[width=.85\linewidth]{aaf.png}
  \caption[Adaptive Acquisition Framework]{Adaptive Acquisition Framework. \srcCite{edo-3-3-1-acq-policy-players-2025}.}
  \label{fig:aaf}
\end{figure}

\subsection{Defense Acquisition System Focus}
The \ac{das} manages the acquisition of weapon systems and other materiel. Major Capability Acquisition is the primary pathway for this course and is event-driven with five phases, four major decision points, and three milestones~\autocite{edo-3-3-1-acq-policy-players-2025}.

\begin{figure}[H]
  \centering
  \includegraphics[width=.75\linewidth]{mca_phases.png}
  \caption[Major Capability Acquisition]{Major Capability Acquisition phases, milestones, and decision points. \srcCite{edo-3-3-1-acq-policy-players-2025}.}
  \label{fig:mca}
\end{figure}

\subsection{Warfighting Acquisition System Emphasis}
\begin{description}
  \item[\textbf{Redesignation.}] The \ac{das} is now referred to as the \ac{was} to reinforce acquisition as a warfighting function focused on rapid capability delivery~\autocite{secwar-was-memo-2025}.
  \item[\textbf{Speed and outcomes.}] Portfolio leaders are expected to trade cost, schedule, and performance for faster fielding when warranted, using only the documentation needed to manage risk and statutory compliance~\autocite{secwar-was-memo-2025,secwar-acq-transformation-2025}.
  \item[\textbf{Commercial-first tools.}] \ac{was} guidance elevates commercial-first sourcing and flexible acquisition tools, including \ac{cso} solicitations and \ac{ot} prototype pathways, to shorten sourcing timelines and widen access to non-traditional vendors~\autocite{secwar-was-memo-2025,DIU-Work-With-Us-2026,USC-4021-OTA}.
\end{description}

\subsection{Key Acquisition Reform Initiatives}
\begin{itemize}
  \item Streamlined \ac{far}/\ac{dfars} updates and continued modernization of the 5000 series.
  \item Mil-Spec review for cost effectiveness and open standards.
  \item Advanced Concept Technology Demonstrations (\ac{actd}) to reduce technical risk.
  \item Control of \ac{toc} in addition to procurement cost.
  \item \ac{ippd} and \ac{ipt} structures to integrate disciplines early.
  \item \ac{caiv} to trade performance and cost.
  \item Single Process Initiative to reduce process duplication.
  \item Open Systems Architecture (\ac{osa}) and simulation-based acquisition to shorten time and cost.
  \item Better Buying Power (\ac{bbp}) and Back to Basics (\ac{btb}) workforce initiatives~\autocite{edo-3-3-1-acq-policy-players-2025}.
\end{itemize}

\subsection{Risk Factor Impact}
\begin{description}
  \item[\textbf{Cost.}] Drivers include material price changes, labor rates, and regulatory compliance such as \ac{esoh} requirements.
  \item[\textbf{Schedule.}] Late deliveries, political pressure, and changing capability needs can drive schedule risk, as can late compliance planning.
  \item[\textbf{Performance.}] New materials or technologies and unproven operational concepts can increase performance risk.
\end{description}
Managing risk by balancing cost, schedule, and performance is central to successful acquisition~\autocite{edo-3-3-1-acq-policy-players-2025}.

\begin{figure}[H]
  \centering
  \includegraphics[width=0.8\linewidth]{cost_schedule_performance.png}
  \caption[Cost, Schedule, Performance]{Cost, schedule, and performance trade space. \srcCite{edo-3-3-1-acq-policy-players-2025}.}
  \label{fig:cost_schedule_performance}
\end{figure}

\subsection{Affordability and Trade Space}
Affordability keeps life-cycle cost within an established range while using schedule and performance as the trade space. It informs \ac{aoa} work, requirements, and engineering decisions throughout the life cycle~\autocite{edo-3-3-1-acq-policy-players-2025}.

\begin{figure}[H]
  \centering
  \includegraphics[width=.8\linewidth]{trade_space.png}
  \caption[Trade Space, Performance vs.\ Cost]{Trade space of performance vs.\ cost showing the knee of the curve. \srcCite{edo-3-3-1-acq-policy-players-2025}.}
  \label{fig:trade_space}
\end{figure}

\subsection{Decision Support Systems}
Three decision support systems must work together for acquisition success~\autocite{edo-3-3-1-acq-policy-players-2025}.
\begin{description}
  \item[\textbf{JCIDS.}] Requirements generation system that identifies, assesses, and prioritizes warfighter capability needs.
  \item[\textbf{PPBE.}] Resource allocation process that builds the \ac{dod} budget for acquisitions.
  \item[\textbf{DAS.}] Management process for developing and buying weapon systems; implemented via the \ac{aaf}.
\end{description}
Program Managers must understand all three systems and synchronize them~\autocite{edo-3-3-1-acq-policy-players-2025}.

\subsection{Major Institutions of Defense Acquisition}
At the national level, defense acquisition is shaped by the Executive Branch, Congress, and defense industry. Program Managers execute day-to-day coordination and reporting across these institutions~\autocite{edo-3-3-1-acq-policy-players-2025}.

\subsection{Defense Industry Landscape}
\begin{description}
  \item[\textbf{Public vs.\ private companies.}] Public companies issue shares and disclose financial results; private companies are closely held and generally have fewer reporting obligations.
  \item[\textbf{Large/medium businesses.}] Establish divisions, team with other large firms, and pursue large \ac{rdte} contracts to gain experience.
  \item[\textbf{Small businesses.}] Often centralize staff functions, outsource support services, and rely on relationships with \acp{ko}, the \ac{sba}, and the \ac{dod} Office of Small Business Programs; the coursebook threshold is 500 employees or \$7.5M in revenue.
  \item[\textbf{Commercial vs.\ government markets.}] Commercial markets set prices freely; government markets are regulated by the \ac{far} and profit caps but are generally less risky~\autocite{edo-3-3-1-acq-policy-players-2025}.
\end{description}

\subsection{Organizational Structures}
Common organizational models include functional, divisional, and matrix structures. Small firms typically centralize functions, while large firms create multiple divisions aligned to product lines or markets. A sound structure clarifies authority and enables coordination; a poor structure blurs accountability and slows decision-making~\autocite{edo-3-3-1-acq-policy-players-2025}.

\subsection{Defense Trends}
\begin{description}
  \item[\textbf{Vulnerabilities.}] Economic downturns, technology breakthroughs, workforce availability, energy issues, raw materials, and wartime demand can stress suppliers and schedules.
  \item[\textbf{Opportunities.}] Shifting threats create opportunities in cyber, counter-terrorism, missile defense, space, and science and technology; creative partnerships can leverage commercial markets~\autocite{edo-3-3-1-acq-policy-players-2025}.
\end{description}

\subsection{Acquisition Players (DoWD 5000.01 Definitions)}
\begin{description}
  \item[\textbf{Acquisition Program.}] A directed, funded effort to deliver a new, improved, or continuing materiel, weapon, information system, or service capability in response to an approved need.
  \item[\textbf{DAE.}] The top-level official for defense acquisitions; acts as \ac{mdauthority} for designated \acp{mdap} and \acp{mais} and may delegate authority to Component leaders.
  \item[\textbf{CAE.}] Senior official responsible for acquisition functions within a \ac{dod} Component; authority may be re-delegated.
  \item[\textbf{MDA.}] Designated official with overall program responsibility; approves phase entry and is accountable for cost, schedule, and performance reporting to higher authority.
  \item[\textbf{PEO.}] Directs multiple acquisition programs and may be delegated \ac{mdauthority} for \ac{acat}~II/III/IV programs.
  \item[\textbf{PM.}] Responsible for achieving program objectives and accountable for cost, schedule, and performance reporting to the \ac{mdauthority}~\autocite{edo-3-3-1-acq-policy-players-2025}.
\end{description}

\subsection{Office of the Secretary of War Decision Makers}
Key decision makers within \ac{osd} include \ac{usdas}, \ac{usdre}, the \ac{dod} \ac{cio}, and \ac{dote}. \ac{cape} provides cost and program evaluation support~\autocite{edo-3-3-1-acq-policy-players-2025}.

\begin{description}
  \item[\textbf{USW(C).}] Controls the budget and release of funds; aligned to \ac{ppbe} budgeting and execution.
  \item[\textbf{USW(P).}] Aligns \ac{dod} plans with national security objectives and supports the planning phase of \ac{ppbe}.
  \item[\textbf{USW(A\&S).}] Principal advisor to \ac{secdef} and Deputy Secretary on acquisition and sustainment; sets workforce policy and makes milestone decisions for designated \acp{mdap}.
  \item[\textbf{DOT\&E.}] Provides independent operational test oversight, approves operational test plans and live fire strategies, and reports directly to \ac{secdef} and Congress.
  \item[\textbf{CAPE.}] Oversees \ac{aoa} guidance and evaluates program costs, schedules, and impacts across the \ac{dod} portfolio.
  \item[\textbf{USW(R\&E).}] Chief technology officer for \ac{dod}; provides technical risk assessments and oversees innovation organizations such as \ac{diu} and \ac{darpa}.
  \item[\textbf{DoW CIO.}] Oversees \ac{it} and cyber acquisition and may serve as \ac{mdauthority} for \ac{mais} programs.
\end{description}
The \ac{dab} is the senior forum advising \ac{usdas} on critical acquisition decisions and includes senior \ac{dod} leadership from policy, financial management, test, intelligence, and Service Secretaries~\autocite{edo-3-3-1-acq-policy-players-2025}.

\subsection{Intelligence Community Touchpoints}
Service and national intelligence organizations provide threat assessments and updated intelligence throughout the acquisition life cycle. Program review involvement typically includes Service intelligence, \ac{dia}, and \ac{dni} depending on program impact~\autocite{edo-3-3-1-acq-policy-players-2025}.

\subsection{Department of the Navy Decision Makers}
\ac{asnrdanda} serves as the \ac{don} \ac{cae} and oversees \ac{syscom} commanders, \acp{peo}, and \acp{dasn}. \ac{asnrdanda} represents the \ac{don} to \ac{usdas} and Congress on acquisition policy and program execution~\autocite{edo-3-3-1-acq-policy-players-2025}.

\subsection{Acquisition Categories}
Programs are grouped into \ac{acat} categories primarily by cost and level of interest to scale oversight and statutory compliance. Factors include cost, location in the acquisition framework, and special interest~\autocite{edo-3-3-1-acq-policy-players-2025}.

\subsubsection{ACAT I}
\ac{acat}~I programs are \acp{mdap}. Current statutory \ac{mdap} dollar thresholds, reset by the \ac{fy}~2026 \ac{ndaa} into FY~2024 constant dollars, are estimated expenditure of at least \$1.0B \ac{rdte} or \$4.5B procurement, or designation by the \ac{usdas}. \ac{acat}~ID programs are managed at the \ac{dae} level; \ac{acat}~IC programs are managed by the \ac{cae}. Examples include DDG-1000 and the Navy Multiband Terminal~\autocite{edo-3-3-1-acq-policy-players-2025,FY2026-NDAA}.

\subsubsection{ACAT II}
\ac{acat}~II major systems do not meet \ac{acat}~I criteria. Current statutory major-system dollar thresholds, reset by the \ac{fy}~2026 \ac{ndaa} into FY~2024 constant dollars, are estimated expenditure of at least \$275M \ac{rdte} or \$1.3B procurement. The \ac{cae} (or designee) is the \ac{mdauthority}. Automated information systems do not use \ac{acat}~II~\autocite{edo-3-3-1-acq-policy-players-2025,FY2026-NDAA}.

\subsubsection{ACAT III and ACAT IV}
\ac{acat}~III programs do not meet \ac{acat}~I, \ac{acat}~IA, or \ac{acat}~II criteria and, in \ac{don} policy, are assigned to the lowest appropriate \ac{mdauthority}, normally the cognizant \ac{peo}, \ac{drpm}, \ac{syscom} commander, or designee. For board answers, flag that coursebook and legacy \ac{don} policy tables may still show FY~2020 ACAT II/III dollar breakpoints, while current statutory major-system screening now uses the FY~2024 \$275M \ac{rdte}/\$1.3B procurement thresholds. \ac{acat}~IIIA covers less-than-major automated information systems. \ac{acat}~IV programs are \ac{don}-unique categories for efforts not designated \ac{acat}~III, and are split into IVT (test) and IVM (monitor) depending on operational test requirements~\autocite{edo-3-3-1-acq-policy-players-2025,SECNAV5000-2G,FY2026-NDAA}.

\subsection{Program Reporting Requirements}
\begin{description}
  \item[\textbf{RDAIS.}] Authoritative \ac{don} database tracking program cost, schedule, and performance with monthly and quarterly updates.
  \item[\textbf{DAES.}] Early warning reports from \ac{don} to \ac{usdas} for \ac{acat}~I programs between milestones, using \ac{rdais} inputs.
  \item[\textbf{DAVE.}] Tracks \ac{mta} efforts with status, issues, cost, and performance updates.
  \item[\textbf{SAR.}] Annual report from \ac{secdef} to Congress for \ac{acat}~I programs; quarterly reporting is required for schedule slips of six months or more or unit cost increases of 15\% over the current \ac{apb} objective (or 30\% over the original baseline).
  \item[\textbf{UCR.}] Quarterly report for \ac{acat}~I programs submitting \ac{sar} to ensure unit costs are visible and not masked by quantity changes~\autocite{edo-3-3-1-acq-policy-players-2025}.
\end{description}

\subsection{Current Navy Acquisition Issues}
\begin{itemize}
  \item China, culture, climate change, and supply chain issues.
  \item ``More than 350'' fleet objective from Force Design 2045 (2022) for manned and unmanned platforms.
  \item Make cybersecurity part of program DNA from requirements through sustainment.
  \item ``Get Real, Get Better'' as a leadership imperative for continuous improvement.
  \item Ongoing acquisition reform~\autocite{edo-3-3-1-acq-policy-players-2025}.
\end{itemize}

\subsection{Summary Highlights}
\begin{itemize}
  \item Decision support systems: \ac{jcids}, \ac{ppbe}, and \ac{das}.
  \item Rules and regulations: \ac{far}, \ac{dodi}~5000.02, and \ac{secnavinst}~5000.2.
  \item Risk focus: cost, schedule, and performance.
  \item Key players: \ac{dae}, \ac{cae}, \ac{mdauthority}, \ac{peo}, and \ac{pm}.
  \item \ac{acat} categories exist to scale oversight and tailor reviews~\autocite{edo-3-3-1-acq-policy-players-2025}.
\end{itemize}

\subsection{Coursebook cue: systems acquisition management}

Coursebook 3.3.1 uses ``Systems Acquisition Management'' to describe the
organized process of translating military need into fielded, supportable
capability through requirements, resourcing, contracting, engineering, test,
production, sustainment, and disposal.  The EDO answer should connect the
management system to three outcomes: deliver required capability, keep the
program executable in cost and schedule, and preserve competitive advantage by
using government contracting and technical authority to buy what the fleet
actually needs rather than what a vendor happens to offer
\autocite{edo-3-3-1-acq-policy-players-2025}.

\paragraph{Board cue.}
Systems acquisition management is the disciplined conversion of need into
fielded capability.  Tie the answer to requirements, resources, contracting,
engineering, test, production, sustainment, and competitive advantage.

%====================
% End of file
